\section{Completion of the proof}
\label{sec:completion}

\subsection{The time-symmetric inequality}

\begin{proof}[Proof of Theorem~\ref{thm:ts}]
The approximation of Proposition~\ref{prop:conformal-approximation}
preserves $H(S)\le0$, converges uniformly in metric comparison, and
preserves the limiting time component. The comparison of enclosing
areas in Section~\ref{sec:preliminaries} therefore allows us first to
assume that $R_g+6>0$ and that the end has the smooth conformal form
used in Section~\ref{sec:stress-extension}. Fix this approximated
metric for the constructions that follow, and put
\[
 A_*=A_{\min}(S;g),\qquad a=\sqrt{A_*/(4\pi)}.
\]

Given an arbitrarily small mass tolerance $\varepsilon_m>0$,
Proposition~\ref{prop:af-attachment} produces a fixed compact inner
piece, a smooth synthetic tensor on it, and a scalar-flat AF
attachment with mass $m_o$ satisfying
\begin{equation}\label{eq:completion-attachment-budget}
 m_o\le p_0(g)-\frac{a^3}{2}+\varepsilon_m.
\end{equation}
The electrostatic height, attachment radius, and smoothing tolerances
used to obtain this inequality have already been fixed. They are
independent of the large graph height in the next construction.

Apply Proposition~\ref{prop:compact-solutions} and the comparison
construction of Proposition~\ref{prop:af-comparison}. Their parameters
are chosen in the proved order: the small inner flux and scalar
cutoff costs first, then a desired relative-area tolerance
$0<\varepsilon_A<1$, a sufficiently small contraction threshold,
a thin omitted exterior collar, and a sufficiently large fixed
volume penalty; next the graph height is taken large, and finally
the lapse clipping, finite penalty strengths, and corner-smoothing
tolerances are chosen. The costs can be made small enough that the
resulting smooth AF comparison metric satisfies
\begin{equation}\label{eq:completion-comparison-budget}
 m_{\AF}(h')\le m_o+\varepsilon_m.
\end{equation}

Let $E$ be its minimizing enclosure containing the caps and the
unsafe-region obstacles. Its free boundary is smooth, compact,
minimal, and outer minimizing. The exterior is complete and has
$R_{h'}\ge0$, the AF decay required in Theorem~\ref{thm:af-penrose},
and one end. Let $E_0$ instead be a largest perimeter-minimizing
enclosure of the caps alone. Then
\begin{equation}\label{eq:completion-enclosure-order}
 \Area_{h'}(\partial E)\ge\Area_{h'}(\partial E_0).
\end{equation}
The exterior of $E_0$ has a scalar lower bound independent of the
large graph height, even where it meets the unsafe region.
The bad-volume estimates from Sections~\ref{sec:compact-construction}
and~\ref{sec:scalar-error}, together with the projection along the
AF attachment, satisfy Proposition~\ref{prop:area-transfer}.
Consequently, when $A_*>0$, the parameters above can be chosen so that
\begin{equation}\label{eq:completion-area}
 \Area_{h'}(\partial E_0)\ge(1-\varepsilon_A)A_*.
\end{equation}
The small initial conformal change used only for the flow fixes this
initial area. It is not made in applying the mass theorem.

Theorem~\ref{thm:af-penrose} applies to the exterior of $E$, not the
possibly unsafe exterior of $E_0$. Equations
\eqref{eq:completion-attachment-budget}--\eqref{eq:completion-area}
give
\begin{align*}
 p_0(g)
 &\ge \frac{a^3}{2}+m_o-\varepsilon_m\\
 &\ge \frac{a^3}{2}+m_{\AF}(h')-2\varepsilon_m\\
 &\ge \frac{a^3}{2}+\frac a2\sqrt{1-\varepsilon_A}
                         -2\varepsilon_m.
\end{align*}
Let the two tolerances tend to zero, making fresh choices of the
finite parameters as necessary. If $A_*=0$, the same construction
and the nonnegativity consequence of Theorem~\ref{thm:af-penrose}
give $p_0(g)\ge-2\varepsilon_m$ without the area step.
Thus \eqref{eq:ts-inequality} holds for every fixed approximated
metric. Passing back through the preliminary approximation proves
it for the original metric.
\end{proof}

\subsection{Maximal data and invariant mass}

\begin{proof}[Proof of the inequality in Theorem~\ref{thm:main}]
The maximal constraint and \eqref{eq:dec-integrability} imply
\[
 \int_{\mathrm{end}} V_0|R_g+6|\dd V_g<\infty.
\]
Use Proposition~\ref{prop:mass-covariance} to choose a rest chart
for the original covector. All hypotheses remain valid in that
chart and $p_0(g)=m_{\AH}(g)$.

Define
\[
 F(A)=\frac12\left(\sqrt{\frac A{4\pi}}
                       +\left(\frac A{4\pi}\right)^{3/2}\right).
\]
This function is continuous, nonnegative, and nondecreasing on
$[0,\infty)$. Theorem~\ref{thm:ts} supplies the time-symmetric input
of Theorem~\ref{thm:ah-reduction}, including comparison metrics whose
covectors are not timelike. The reduction therefore gives
\[
 p_0(g)\ge F\bigl(A_{\min}(S;g)\bigr).
\]
Since the chart was fixed before any construction, this is precisely
\eqref{eq:main-inequality}. The one-sided time-component budgets
suffice; there is no further assertion about the spatial components
of intermediate comparison metrics.
\end{proof}

\subsection{Sharpness}

Let $m>0$, and let $a$ be the positive solution of
$a+a^3=2m$. The Schwarzschild--anti-de Sitter metric is
\begin{equation}\label{eq:sads-metric}
 g_m=\frac{\dd r^2}{1+r^2-2m/r}+r^2\sigma,
 \qquad r\ge a,\qquad K=0.
\end{equation}
To include its boundary smoothly, use proper distance $s$ from
$r=a$. Then
\[
 g_m=\dd s^2+r(s)^2\sigma,
 \quad r(0)=a,\quad r'(0)=0,
 \quad (r')^2=1+r^2-2m/r.
\]
The simple positive radial derivative of the expression on the
right at $a$ gives a smooth even solution $r(s)$ near zero, increasing
for $s>0$. Hence the horizon is a smooth minimal boundary and the
exterior is complete with that boundary included. Since
$r''=r+m/r^2$, the warped-product formula gives
\[
 R_{g_m}=-\frac{4r''}{r}
             +\frac{2(1-(r')^2)}{r^2}=-6.
\]
The constraints are vacuum, the metric error and its required tensor
derivatives are $O(r^{-3})$, and the weighted matter integral vanishes.

For completeness the mass normalization can be checked directly.
In the background distance coordinate, write
$g_m-b=\zeta\dd\eta^2$, where
\[
 \zeta=\frac{2m}{r(1+r^2-2m/r)}.
\]
The radial component of $\divg_b e-\dd\tr_be$ is
$2\coth\eta\,\zeta$, and the other two terms of
\eqref{eq:mass-integrand} cancel for $V_0$.
The normalized sphere flux therefore tends to
\[
 m\frac{1+r^2}{1+r^2-2m/r}\longrightarrow m.
\]
Rotational symmetry makes the three spatial fluxes zero, so
$m_{\AH}(g_m)=m$.

Radial projection onto the horizon has two-area Jacobian at most one,
since $r\ge a$. Every radial ray meets the entire intrinsic boundary
of any admissible end-side domain: either it enters that domain
across its boundary, or its initial point already lies on the counted
boundary at $S$. Thus projection of the cut covers the horizon.
The area formula gives $\Area_{g_m}(\partial D)\ge4\pi a^2$.
The horizon itself is an admissible competitor, and hence
$A_{\min}(S;g_m)=4\pi a^2$. Equation $2m=a+a^3$ proves equality in
\eqref{eq:main-inequality} and completes Theorem~\ref{thm:main}.
